\section{Deletion Features in E2EE Messaging}\label{sec:e2ee-deletion-systematization}
We begin with a broad overview of message deletion in E2EE applications, by providing a \emph{systematization of message deletion features} that distills their core properties and differences. Looking ahead, the taxonomy we introduce in this section will serve to motivate our forthcoming threat model.

\paragraph{What are message deletion features?} Before presenting our systematization, we first define what we mean by a message deletion feature. While seemingly trivial, not all features that hide messages from the user interface fall under the umbrella of ``deletion'', so we first specify the criteria for what is in scope of our work.

At a high level, we define a \emph{message deletion feature} as any application functionality\footnote{By ``application functionality'', we refer to the actions exposed to users via the client's user interface.} that is \emph{intended} to make messages \emph{unrecoverable}. By ``unrecoverable'', we mean that \emph{no other application-level feature allows for retrieving messages once they have been deleted}. Notice that our definition is based on what a feature is \emph{designed} to provide (hence our usage of the term ``intended''); whether messages are \emph{actually} recoverable or not depends on the (security of the) specific implementation. (Formally defining the \emph{security guarantees} of deletion features is the focus of Section~\ref{sec:security-definition}.)

For example, the most common message deletion feature is \emph{chat deletion},\footnote{In this work, we use ``chat'' and ``conversation'' interchangeably to refer to the list of all individual messages exchanged between two or more users.} which erases all messages in a conversation and removes the conversation from the list of active chats. In addition, many applications (e.g., WhatsApp and Signal) provide more fine-grained deletion features that allow for removal of individual messages. These include \emph{local deletion}, which allows users to remove any (sent or received) message from their view of the conversation; as well as \emph{unsending}, which allows the sender of a message to remove it for both them and the recipient(s), typically replacing it with a placeholder string (e.g., ``this message was deleted'' in WhatsApp~\cite{whatsapp-delete-message}). 

Furthermore, while chat deletions, local message deletions, and unsending are triggered manually by users, many applications additionally provide features for \emph{automatic} deletion of messages. These include \emph{disappearing messages}, which allow users to specify a lifetime for messages, after which they are removed from the chat for all parties. Another example is \emph{view-once media}, which allows users to send media attachments (pictures, videos, etc.) that can only be opened once by recipients.

% (In Appendix~\ref{}, we provide detailed context of how these features work, including walkthroughs of how to use them in two case-study applications.).


As mentioned earlier, whether a feature constitutes deletion or not is relative to its intended behavior, and so it is ultimately application-specific. Often, features have consistent goals across applications; for instance, as far as we are aware, no E2EE application supports recovering messages for any of the features mentioned above. For example, WhatsApp explicitly states that ``when you delete a chat, it can't be undone''~\cite{whatsapp-delete-chat}, and that when a message is unsent ``there's no way to get it back''~\cite{whatsapp-delete-message}. However, in other cases, features may have different goals in different applications. For example, many applications allow users to \emph{edit} messages; in some applications, users can read the history of edits (e.g., in Signal~\cite{signal-edit-message}), while others only display the latest version (e.g., in WhatsApp~\cite{whatsapp-edit-message}). So, edits would be a deletion feature in the latter but not the former.

On the other hand, our definition of message deletion excludes features such as, for example, \emph{archiving} conversations: while this hides conversations from the list of active chats, users can unarchive (and, thus, reinstate) them later on. Indeed, WhatsApp explicitly clarifies that ``archiving a chat doesn't delete the chat''~\cite{whatsapp-archive-chat}. As a second non-example, features for hiding chats in a private (often password-locked) section of the UI (e.g., WhatsApp's ``locked chats'' feature~\cite{whatsapp-locked-chat}) are also not in scope of our definition.

\paragraph{Discussion.} Our criteria for message deletion has a few subtleties, which will be relevant to our threat model in Section~\ref{sec:threat-model}. First, we focus on the behavior of features directly on \emph{original} messages, and not on \emph{derived} content. For example, \emph{replies} are a common feature in messaging applications, which typically embed the first characters of the original message alongside the new message. We treat replies and other derived content as logically different from the original message, and thus do not define a message deletion feature in terms of its ability to remove message references in replies.

Second, we define message deletion in terms of the \emph{local, present state} of clients. For any feature, deleted messages may of course be recoverable if the user has an older copy of all their messages (e.g., a cloud backup) from before deletion occurred. We do not deem this as reverting deletions, but rather restoring an old device state. Even if subsequent backups do not include deleted messages, users can manually store older backup versions, which are outside the reach of deletion features.


Similarly, for simplicity, we focus on the \emph{single-device setting}.



% A paragraph with remarks, either here or at the end of section.
% Interplay with message replies/quotes.
% Interplay of our definition with backups.
% Interplay with screenshots


\paragraph{Systematization of message deletion features.} Having defined message deletion features, we now turn to their properties; while deletion features share high-level goals, their specific functioning varies. So, we now introduce an \emph{analytical framework} that systematizes the key properties of deletion features. Our goal is to highlight the main design choices and differences among such features, to facilitate their systematic understanding.


In a nutshell, we identify four main properties of message deletion features: (1) which users can trigger the feature, (2) the endpoint(s) where deletion take place, (3) whether deletion is silent or visible, and (4) how long after feature activation deletion takes place. So, together, these properties logically capture the ``who'', ``where'', ``what'', and ``when'' of a deletion feature. We describe each property in turn below, and show example deletion features and how they map onto our systematization in Table~\ref{fig:deletion-features-systematization-examples}. While, as mentioned earlier, the specific implementation of a feature (and, thus, how it fits within our framework) may vary across applications, we analyze them in the context of our two-case studies, WhatsApp and Signal, which are representative of the ecosystem more broadly.

\begin{figure*}[t!]
    \centering
    % \renewcommand{\arraystretch}{1.2}
    \begin{tabular}{l|lllll}
     \toprule
     \textbf{Feature} & \textbf{Initiating endpoint} & \textbf{Scope of removal} & \textbf{Replacement} & \textbf{Timing} & \textbf{P5} \\ 
     \midrule
    Close conversation & Sender and recipient-side & A & A & A & A \\
    Unsend message & Sender-side & Local & A & A & A \\
    Delete message & Sender and recipient-side & Global & A & A & A \\
    Timed message & A & A & A & A & A \\
    \bottomrule
\end{tabular}
\caption{Example of common message deletion features, and how they fit within our systematization of Section~\ref{sec:e2ee-deletion-systematization}.}\label{fig:deletion-features-systematization-examples}
\end{figure*}

\smallskip\emph{Property \#1: initiating endpoint.} The first property defines which endpoint of a message can trigger the deletion feature. A feature is \emph{sender-based} if only the sender of a message can trigger deletion (e.g., unsending a message); it is \emph{recipient-based} if an action by the recipient triggers deletion (e.g., view-once media); or it is \emph{endpoint-independent} if it can be triggered by any party in the conversation (e.g., local message deletion).

\smallskip\emph{Property \#2: deletion scope.} The second property defines whether the feature propagates deletion beyond the initiating endpoint. A feature has \emph{local} scope if it only modifies the view of the user that triggers it (e.g., chat deletion); or it has \emph{full} scope if it modifies the view of all users in the conversation (e.g., disappearing messages).


\smallskip\emph{Property \#3: deletion visibility.} The third property defines the way in which usage of the feature is reflected in the user interface. A feature is \emph{silent} if it removes messages without any visual indication, as if the message had never been sent (e.g., local message deletion); or it is \emph{visible} if it replaces messages with a placeholder string, indicating that deletion ocurred but without revealing the content of the message (e.g., unsending).

\smallskip\emph{Property \#4: time interval.} Lastly, the fourth property specifies the \emph{amount of time} between when the feature is triggered and when messages are deleted. A feature is \emph{instantaneous} if messages are deleted the moment the feature is triggered (e.g., unsending); or it is \emph{timed} if it retains messages for a defined period of time before they are removed (e.g., timed disappearing messages). \smallskip

In addition to these four properties, there are other behaviors that may vary across message deletion features. For example, features may differ in the amount of time in which they can be invoked after a message is sent (e.g., for unsending, 24 hours on Signal~\cite{signal-delete-message} and 48 hours on WhatsApp~\cite{whatsapp-delete-message}), or in whether they operate on bulk or individual messages (e.g., chat deletion vs local message deletion), etc. However, we focus on the four properties above as they are the most directly relevant to reasoning about deletion security.